\section{The monodromy obstruction and its integral inverse}

% BEGIN M85N1
\paragraph{M85N1. The induced Euler operator preserves the integral principal parts.}
Use the realized Breuil--Kisin data of Q84T2--Q84T4. Write
$v=u^e$, $E=v+p$, $A=O[[u]]$, $L=A\varphi M$ and
$\mathcal L=SL$, with the unramified components understood.
Here $S$ is the Breuil ring and $R_\pi$ is the analytic ring
of the closed disc of radius $p^{-1/e}$ over the coefficient field.
In the Frobenius-column basis $\mathcal L$ is free over $S$.
Write $M_{\rm an}=R_\pi M$, $L_{\rm an}=R_\pi L$,
and $[x]$ for the class modulo $M_{\rm an}$.
On these principal parts $T[x]=[E\nabla x]$, where
$\nabla$ has scalar derivation $u\,d/du$.
Put
\[
 Y=(E^{-b}M)^\Delta,\quad Y'=vY,\quad
 Z=\bigl([E^{-b}M]\bigr)^\Delta,\quad Z_L=Z[1/p].
\]
The order of the finite tame descent group $\Delta$ is a $p$-unit
in Q84T2 (both $e$ and the residue degree $a$ are $p$-units).
Thus averaging gives $[Y]=Z$.

The series $f=\sum f_nu^n$ in $S$ have
$f_n\lfloor n/e\rfloor!\in O$ tending to zero.
Write $n=eq+r$, $0\le r<e$, and expand $u^n=u^r(E-p)^q$.
In the coefficient of $u^r E^i/i!$ the multiplier of $f_n q!$
is $(-p)^{q-i}/(q-i)!$ for $i\le q$, hence integral.
The sums converge and their tails tend uniformly to zero by
the bound on $f_nq!$. For $j<p$, the terms with
$i<j$ form an integral polynomial in $u$ because $i!$ is a unit.
The other terms belong to the closed divided-power filtration
$\mathrm{Fil}^jS$. In $R_\pi$ they are multiples of $E^j$:
$S$ maps continuously to the disc algebra by the factorial bound,
and the analytic ideal $(E^j)$ is closed by monic division on
that disc, equivalently by the finitely many root jets after
splitting $E$. Hence
\[
 \mathcal L=L+\mathrm{Fil}^{10H}\mathcal L,\qquad
 E^{-2H}\mathrm{Fil}^{10H}\mathcal L\subset M_{\rm an}.
\]
The second inclusion uses $E^b L\subset M$ and $8H\ge b$.
Also $p^b\varphi\mathcal L\subset\mathcal L$:
$E^b\mathcal L\subset SM$, $\varphi(S)\subset S$ by the factorial
bounds, and $\varphi(E)/p=1+u^{ep}/p$ is a unit in $S$.
Indeed the $n$th power of $u^{ep}/p$ has weighted coefficient
of valuation $v_p((pn)!)-n=v_p(n!)$, tending to infinity.

For $y\in E^{-b}M$ let $m=E^{H+b}y\in L$ by $E^H M\subset L$.
The Leibniz rule gives
\[
 E^{2H}(E\nabla y)
   =E^{H-b+1}\nabla m-(H+b)E^{H-b}(uE')m .
\]
The first term belongs to $\mathcal L$ by the connection condition
of the realized model, and the second belongs to $L$.
Split the result as $l+f$ with $l\in L$ and
$f\in\mathrm{Fil}^{10H}\mathcal L$. The term $E^{-2H}f$ is
in $M_{\rm an}$. Also the model gives
$E\nabla M_{\rm an}\subset M_{\rm an}$, hence
$E\nabla(E^{-b}M_{\rm an})\subset E^{-b}M_{\rm an}$ by Leibniz.
Thus $E^bE^{-2H}l\in M_{\rm an}$. This vector also belongs to
$E^{-2H}M$ since $E^bL\subset M$. Their intersection is $M$:
divide the integral numerator by the distinguished polynomial
$E^{2H}$ in each coordinate. Analytic divisibility forces the
polynomial remainder to vanish at every root with the prescribed
multiplicity. We have shown $T[y]=[E^{-2H}l]\in [E^{-b}M]$.
Since $E\nabla$ commutes with descent, $TZ\subset Z$.
% END M85N1

% BEGIN M85N2
\paragraph{M85N2. The shifted jump spectrum on the rational principal parts.}
Define $D_U=\{i+k:i\in U,\ 1\le k\le b\}$, taking each
integer just once. On $Z_L$ the operator $T$ is semisimple,
with eigenvalues among $nep$ for $n\in D_U$.
To see this put $W=M_{\rm an}$,
$W'=E^{-b}W$, $F^j=E^jL_{\rm an}$. On $F^j/F^{j+1}$ the
operator $E\nabla$ acts by $-jep$. Indeed $uE'=e(E-p)$, so
\[
 (E\nabla+jep)(E^j x)=E^{j+1}(je x+\nabla x).
\]
The analytic Frobenius-column lattice is stable under $\nabla$:
the Euler derivative preserves convergent series on the disc,
and the derivatives of the basis vectors belong to $\mathcal L$.
The height inclusions give $F^b\subset W\subset W'\subset
F^{-H-b}$.

On $L_{\rm an}/E L_{\rm an}$ the filtration given by
$L_{\rm an}\cap E^i M_{\rm an}$ is the original jump
filtration evaluated at $E=0$. To check this, divide the $M_{\rm an}$
coordinates of $E^{-i}x$ by $E^h$, $h=\max(0,H+1-i)$, for
$x\in L_{\rm an}\cap E^iM_{\rm an}$. Multiplying the polynomial
remainders by $E^i$ gives a vector agreeing with $x$ modulo
$E L_{\rm an}$ because $E^H M\subset L$. This approximation
belongs to $L[1/p]$: its Frobenius-column coordinates have a power
of $E$ as their only possible denominator over $A[1/p]$, and
analytic divisibility removes it by the remainder argument in M85N1
after clearing coefficient denominators. It also belongs to
$E^iM[1/p]$. The reverse inclusion follows by extending scalars.

For $x\in W'\cap F^j$, the residue of $E^{-j}x$ lies in
filtration degree $-j-b$. When there is no jump in
$[-j-b,-j-1]$, that filtration equals its degree $-j$.
Subtracting a representative in $W\cap F^j$ then puts $x$ in
$F^{j+1}$. When a jump occurs, $-j\in D_U$, and the factor
$(E\nabla+jep)$ puts $x$ in $F^{j+1}$ modulo $W$.
All factors preserve $W,W'$. Moving successively from
$j=-H-b$ to $b$ shows that
$\prod_{n\in D_U}(T-nep)=0$ on $W'/W$, hence on the invariant
principal parts $Z_L$. Its roots are distinct in characteristic
zero, proving the assertion. For a $T$-stable rational
subspace $V$ one has $V\cap\mathrm{im}\,T=T(V)$: the polynomial
projectors onto the nonzero eigenspaces and the inverse on each
such eigenspace preserve $V$.
% END M85N2

% BEGIN M85N3
\paragraph{M85N3. An integral inverse for the leading coefficient.}
Put $h_L=(N+1)b$, $R(X)=\prod_{0\ne n\in D_U}(X-nep)$ and
$R_1=(R-R(0))/X$. For nonzero $n$ in the product,
$|n|\le H+b<p$ and $p\nmid e$. Thus
$v_p(R(0))=|D_U\setminus\{0\}|\le b|U|\le h_L$,
also for the empty product. If $z=Tw\in Z_L$ with $p^c z\in Z$,
M85N2 gives $R(T)z=0$. Consequently
\[
    T\bigl(-R_1(T)p^c z\bigr)=R(0)p^c z .
\]
The vector inside $T$ is integral by M85N1. Multiplying
it by the integral scalar $p^{h_L}/R(0)$ gives $z_1\in Z$
with $Tz_1=p^{c+h_L}z$. Choose $y_1\in Y$ with $[y_1]=z_1$.
Since $E^b y_1\in M$, the identity
\[
  p^b=E^b-v\sum_{i=0}^{b-1}p^i E^{b-1-i}
\]
shows that $p^b z_1$ is represented by an element of $Y'=vY$.
Thus for each $c\ge0$ and $w\in Z_L$ such that $p^c Tw\in Z$,
some $y\in Y'$ satisfies $T[y]=p^{c+b+h_L}Tw$.
% END M85N3

% BEGIN M85N4
\paragraph{M85N4. The error series on the shifted invariant lattice.}
Set $k=p-14H$ and
$\zeta y=-\sum_{n\ge0}(p\varphi)^{n+1}\nabla y$.
This is the native formal series: $\nabla$ has zero constant
coefficient on every vector and $\varphi$ multiplies its order at
$u=0$ by at least $p$. In particular the sum is coefficientwise
well defined. If $y=v y_*$ with $E^b y_*\in M$, then
\[
       \zeta y\in p^k L+p^{-3H}\mathrm{Fil}^{10H}\mathcal L .
\]
For $w_n=(p\varphi)^{n+1}\nabla(vy_*)$ put
$q_n=H+2b+n(b-1)$. Then
\[
            p^{q_n}w_n\in u^{ep^{n+1}}\mathcal L .
\]
Indeed $E^{H+b}y_*\in L$. Differentiating this equality in
$\mathcal L$ by Leibniz gives
$E^{H+b+1}\nabla y_*\in\mathcal L$, hence
$E^{H+2b+1}\nabla y_*\in SM$. Also $E^b y_*\in M$.
As $\varphi(E)/p$ is a unit of $S$ and $\varphi(SM)\subset
\mathcal L$, it follows that $p^{H+2b+1}\varphi(ey_*+\nabla y_*)$
is in $\mathcal L$. Now $\nabla(vy_*)=v(ey_*+\nabla y_*)$,
giving the formula at $n=0$. Apply $p\varphi$ and
$p^b\varphi\mathcal L\subset\mathcal L$ for the induction.

The factorial bound on coefficients gives
$u^{ep^{n+2}}\in p^{p^{n+1}} S$, since
$v_p((p^{n+2})!)\ge p^{n+1}$. Therefore
$w_{n+1}\in p^{k+n}\mathcal L$ for $n\ge0$. Indeed
$q_{n+1}+k+n\le (n+1)p\le p^{n+1}$, using $b<H$ and
$14H<p$. These terms sum in $p^k\mathcal L$:
the weighted coefficient sequences of $S$ form the complete
space of null sequences with entries in $O$, with its supremum
norm. Convergence there gives the same formal limit coefficient
by coefficient. For $w_0$, expand
$u^{ep}=(E-p)^p$. Terms of $E$-degree at least $10H$ lie in
$\mathrm{Fil}^{10H}S$ and the others are divisible by $p^{p-10H}$.
Hence
\[
 w_0\in p^{p-10H-H-2b}\mathcal L+
                p^{-H-2b}\mathrm{Fil}^{10H}\mathcal L.
\]
Since $H+2b\le3H$ and $p-11H-2b\ge k$,
splitting $\mathcal L=L+\mathrm{Fil}^{10H}\mathcal L$ by M85N1
proves the assertion, including the sign in $\zeta$.
In particular $[E\zeta y]\in p^k [E^{-b}M]$: multiplying
the filtration term by $E$ gives zero modulo $M_{\rm an}$,
and $EL\subset E^{-b}M$. For $y\in Y'$ this class divided
by $p^k$ is invariant by equivariance of the operators.
% END M85N4

% BEGIN M85N5
\paragraph{M85N5. The tame orbit bound for the Frobenius term.}
Write $R$ for the maximal Frobenius period of the character
differences in Q84T2, and $c_1=b(R+1)$. The zero character
belongs to the set, so every occurring inertia character has
period at most $R$. For all $m\in M^\Delta$ we claim
\[
                   p^{c_1}\zeta(\varphi m)\in\mathcal L .
\]
Let $x=\nabla(\varphi m)$. This belongs to $\mathcal L$, is
inertia-invariant and has positive order at $u=0$. If $x=0$
there is nothing to prove. If its least order in the components
is $j<e$, the nonzero leading coefficient of $u^{-j}x$
has tame character $-j\bmod e$ on the fibre at zero: inertia
scales $u$ by its defining root, and reduction at zero commutes
with its action. For the period $r\le R$ of this character,
$e\mid j(p^r-1)$, so $j p^R\ge e$. If $j\ge e$ use $s=e$,
otherwise $s=j$. We have $1\le s\le e$, $e\le sp^R$, and
$x$ is divisible by $u^s$ in formal coordinates.

Write $x$ in the Frobenius-column basis, which is invertible over
the coefficient field's formal power series since $E(0)=p\ne0$.
Its coordinates $f\in S$ are divisible by $u^s$ there. The
quotient $f(u^p)/u^{sp}$ belongs to $S$. In fact at degree
$pn$ its coefficient is $f_{n+s}$; the weighted coefficient
bound follows from
$v_p(\lfloor pn/e\rfloor !)\ge v_p(\lfloor(n+s)/e\rfloor!)$.
For $(p-1)n\ge s$ this follows from monotonicity; otherwise
$n+s<2e$ (also for $p=2$), and the latter factorial is a unit.
The tail condition follows from that for $f$. Applying Frobenius
to its basis vectors costs at most $p^b$ by M85N1. This gives
$\varphi x=u^{sp}x_1$, $p^b x_1\in\mathcal L$.
Consequently
\[
 (p\varphi)^{n+1}x
   =p^{n+1}u^{sp^{n+1}}\varphi^n x_1,\qquad
 p^{(n+1)b}\varphi^n x_1\in\mathcal L .
\]
For $n\le R$, $p^{c_1}$ clears this denominator. For
$n=R+t$, $t\ge1$, the variable power has exponent at least
$ep^{t+1}$. Its factorial weight supplies at least $p^t$
additional powers of $p$. After multiplying by $p^{c_1}$,
the net exponent is at least
$R+t+1+p^t-tb\ge t$, since $p^t\ge tp\ge tb$.
This proves convergence in $\mathcal L$ and the claimed
bound on the full error vector.

The native obstruction satisfies
$\mathrm{Mon}(z)=T[z]-[E\zeta z]$. For $y\in Y'$,
M85N1 and M85N4 give $\mathrm{Mon}(y)\in Z$.
For $m\in M^\Delta$, $\varphi m\in Y$ by the upper height
bound, so $T[\varphi m]\in Z$ too. The calculation above
puts $p^{c_1}E\zeta(\varphi m)$ in $E\mathcal L$, whose
principal parts come from $EL$ by M85N1. Since $EL\subset
E^{-b}M$, taking invariants yields for every such pair
\[
      p^{c_1}\mathrm{Mon}(y+\varphi m)\in Z .
\]
% END M85N5

% BEGIN M85N6
\paragraph{M85N6. Lifting the leading coefficient to exact monodromy.}
For the realized datum let
$W_0=Z\cap T(Z_L)$, a $p$-saturated submodule of $Z$:
the second factor is a vector subspace. K85F4 gives the rational
monodromy-range identity $\mathrm{Mon}(Y[1/p])\subset T(Z_L)$,
and K85F1 gives $p$-adic completeness of $Y'$ and separatedness of $Z$.
Put $A_0 y=T[y]$. By M85N1 and M85N4 both $A_0$ and
$B_0=p^{-k}(\mathrm{Mon}-A_0)$ carry $Y'$ into $Z$.
The range identity and saturation give $B_0(Y')\subset W_0$.
M85N3 at $c=0$ gives $p^{d}W_0\subset A_0(Y')$ for
$d=b+h_L=(N+2)b$. The native prime guard gives $k>d$:
indeed $14H+(N+2)b<p$ follows from $b<H$, $N\ge2$ and
$40(HN^2+N^3)<p$.

For any $y\in Y'$ start with $x=y$, aiming at
$\mathrm{Mon}(x)=A_0 y$. The error is initially in $p^kW_0$.
When the remaining error is $p^m z$ with $z\in W_0$ and $m\ge k$,
choose $a\in Y'$ with $A_0a=p^d z$. Add $p^{m-d}a$ to the
approximation, taking the error as target minus current value.
The new error belongs to $p^{m+k-d}W_0$ because
$\mathrm{Mon}=A_0+p^kB_0$. Each correction thus increases
its guaranteed precision by $k-d>0$. The corrections converge
in $Y'$. Both integral maps are continuous for the $p$-adic
topologies, and separatedness gives the asserted equality in $Z$.
Combining this exact correction with M85N3, for every
$w\in Z_L$ and $c\ge0$ with $p^cTw\in Z$ there exists
$y_0\in Y'$ satisfying
\[
                \mathrm{Mon}(y_0)=p^{c+b+h_L}Tw .
\]
The argument works on the kernel datum after restricting its jumps
to $U$. Its cost beyond the clearing exponent $c$ is $b+h_L$.
% END M85N6
